\section{Stock Span Problem}
\label{sec:sq-stock-span}

\noindent We now turn to the final group of problems in this chapter: a
handful of practical implementation and design problems that each combine
stacks or queues with a hash map or other auxiliary structure to solve a
real-world-flavored task efficiently.

\problemstatement{We are given the price of a stock on successive days,
arriving one at a time (a stream). For each new day's price, compute its
\textbf{span}: the number of consecutive days, ending with and including
today, for which today's price is greater than or equal to each of those
days' prices.}

\begin{patternbox}
\emph{``For each new arrival, look backward and count a contiguous run
based on a comparison''} is the \emph{previous} greater/smaller-element
family from Section~\ref{sec:sq-next-smaller}, but with a twist: we need
a \textbf{count} (the span length), not just the nearest qualifying
index. The key efficiency idea is that, instead of storing every
individual day on the stack, we can store \textbf{(price, span) pairs}
and merge spans together as we pop -- avoiding ever re-scanning days that
have already been absorbed into a previous day's span.
\end{patternbox}

\subsection*{Understanding the problem carefully}

The span of today's price is $1$ (counting today itself) plus the spans
of every immediately preceding day whose price is $\le$ today's price,
chained backward as far as this holds. If we stored only raw prices on
the stack, computing this span naively would require popping one day at a
time and could cost $O(n)$ per query in the worst case, but with no way
to skip over runs that were already resolved. Storing the \emph{span}
\emph{alongside} each price lets us absorb an entire previously-resolved
run in a single pop, rather than one day at a time.

\begin{ideabox}[Stack of (Price, Span) Pairs, Merging on Pop]
Maintain a stack of $(\textit{price}, \textit{span})$ pairs. For each new
price $p$: initialize $\textit{currentSpan} \gets 1$. While the stack is
not empty and its top's price $\le p$: pop the top, and add its
\textit{span} to $\textit{currentSpan}$ (this whole already-resolved run
is absorbed in one step, not re-examined day by day). Push
$(p, \textit{currentSpan})$, and report $\textit{currentSpan}$ as today's
span.
\end{ideabox}

\subsection*{Why absorbing a popped pair's entire span (not just its price) is correct}

If the top of the stack is a pair $(\textit{price}', \textit{span}')$
with $\textit{price}' \le p$, then \emph{every} day within that pair's
already-computed span also has a price $\le \textit{price}' \le p$ (that
is precisely what the pair's span certifies, by the inductive
construction of the algorithm) -- so all of those days extend today's
span too, and we can add the whole count $\textit{span}'$ in one step
rather than re-deriving it. This amortizes the total work: each day is
pushed exactly once and, across its entire lifetime, can only be merged
away (absorbed into a later pair) once, giving linear total work despite
some individual days having large spans.

\subsection*{Algorithm}

\begin{enumerate}
  \item Initialize an empty stack of (price, span) pairs.
  \item For each new price $p$:
  \begin{enumerate}
    \item $\textit{currentSpan} \gets 1$.
    \item While the stack is not empty and its top's price $\le p$: pop
          $(\textit{price}', \textit{span}')$;
          $\textit{currentSpan} \gets \textit{currentSpan} +
          \textit{span}'$.
    \item Push $(p, \textit{currentSpan})$.
    \item Report $\textit{currentSpan}$.
  \end{enumerate}
\end{enumerate}

\subsection*{Dry run}

\dryrun{Prices arrive in order: $100, 80, 60, 70, 60, 75, 85$.}

\begin{itemize}
  \item $100$: stack empty. span$=1$. Push $(100,1)$. Report $1$.
  \item $80$: top $(100,1)$, price $100>80$, no merge. span$=1$. Push
        $(80,1)$. Report $1$.
  \item $60$: top $(80,1)$, $80>60$, no merge. span$=1$. Push $(60,1)$.
        Report $1$.
  \item $70$: top $(60,1)$, $60\le70$: merge, span$=1+1=2$, pop. New top
        $(80,1)$, $80>70$, stop. Push $(70,2)$. Report $2$.
  \item $60$: top $(70,2)$, $70>60$, no merge. span$=1$. Push $(60,1)$.
        Report $1$.
  \item $75$: top $(60,1)$, $60\le75$: merge, span$=1+1=2$, pop. Top
        $(70,2)$, $70\le75$: merge, span$=2+2=4$, pop. Top $(80,1)$,
        $80>75$, stop. Push $(75,4)$. Report $4$.
  \item $85$: top $(75,4)$, $75\le85$: merge, span$=1+4=5$, pop. Top
        $(80,1)$, $80\le85$: merge, span$=5+1=6$, pop. Top $(100,1)$,
        $100>85$, stop. Push $(85,6)$. Report $6$.
\end{itemize}

Spans: $[1,1,1,2,1,4,6]$.

\subsection*{C++ implementation}

\begin{lstlisting}
#include <bits/stdc++.h>
using namespace std;

class StockSpanner {
    stack<pair<int,int>> st; // (price, span)

public:
    int next(int price) {
        int span = 1;
        while (!st.empty() && st.top().first <= price) {
            span += st.top().second;
            st.pop();
        }
        st.push({price, span});
        return span;
    }
};

int main() {
    StockSpanner spanner;
    for (int p : {100, 80, 60, 70, 60, 75, 85}) {
        cout << spanner.next(p) << " ";
    }
    cout << endl;
    return 0;
}
\end{lstlisting}

\begin{complexitybox}
\textbf{Time Complexity.} Each day is pushed once and merged away
(popped) at most once across its entire lifetime, giving
\[
  O(1) \text{ amortized per day}, \qquad O(n) \text{ total for } n \text{
  days}.
\]
\textbf{Space Complexity.} $O(n)$ for the stack in the worst case
(strictly decreasing prices, where every day stays on the stack).
\end{complexitybox}

\begin{takeawaybox}
When a monotonic-stack problem needs a \emph{count} of absorbed elements
rather than just the nearest boundary, store that count directly inside
each stack entry and merge counts on pop -- this avoids re-deriving the
count from scratch every time and keeps the whole algorithm at amortized
$O(1)$ per arrival, even for a problem (like this one) where the answer is
a span length, not a position.
\end{takeawaybox}
